% !PRINTED 122
\weberpage{122}
%<<CONT
prove here by a method of \gesperrt{Frobenius} which rests only on the general concept of a group\footnote{\gesperrt{Cauchy}, Exerc. d'analyse, Tom. III. \gesperrt{Sylow}, Mathem. Annalen, Bd. 5. \gesperrt{Frobenius}, Journ. für Mathematik, Bd. 100. \gesperrt{Netto}, Mathem. Annalen, Bd. 13; Substitutionentheorie. \S.\,48.}.

The theorem can be stated simply thus:

\gesperrt{I. If $P$ is a group of any elements whatever of degree $n$, and $p^{\varkappa}$ is a prime power dividing $n$, then the group $P$ has a divisor of degree $p^{\varkappa}$.}

\gesperrt{Cauchy} has proved this theorem for $\varkappa = 1$. For the case that $n$ is a prime, it is evident; for $n = 4$ and $n = 6$ one can read it off easily from the compilations in \S.\,27, so that it may therefore be regarded as proved for the cases $n = 2, 3, 4, 5, 6, 7$. On the basis of this observation, complete induction can be applied, and we therefore assume \gesperrt{that the theorem has been proved for groups whose degree is lower than} $n$.

We shall denote the elements of the group $P$ by the letters $a, b, c, \ldots$ and understand by $x$ a symbol for a variable element in $P$ which runs through the whole group $P$.

We first gather together all those elements $a$ which satisfy the condition
\weq{1}{x^{-1}ax = a \quad \text{or} \quad ax = xa}
so that $a$ denotes every element that is \gesperrt{permutable} with all elements of $P$. To these belongs in any case the unit element, and the totality of these elements $a$ forms a group $A$ contained in $P$. For from
\[ x^{-1}ax = a,\quad x^{-1}a'x = a' \]
it follows that
\[ x^{-1}aa'x = aa'. \]

\gesperrt{$A$ is an Abelian group and a normal divisor of $P$}; for from the definition (1) it follows that $x^{-1}Ax = A$, and for any two elements $a, a'$ of $A$:
\[ aa' = a'a. \]
We denote by $\nu$ and $j$ the degree and index of this group, so that
\weq{2}{n = \nu j}
%CONT>>
